%% kapa: exam=2 round=24 year=2013 period=3 subject=compensation_law subject_ko=감정평가 및 보상법규 seq=2 number=2 points=30
%% source: Q-Net 감정평가사 기출문제 (cst00309) artlSeq=5077404 file=extracted/1008/5077404/2013년 제24회 감정평가사 2차시험 기출문제/3교시.hwp
%% conversion: kordoc 4.24.51 -> tools/gen2.py (auto); review=proofread-sonnet; flags=none
\begin{longquestion}{【문제 2】}
甲은 S시에 임야 30,000m\textsuperscript{2}를 소유하고 있다. S시장은 甲 소유의 토지에 대하여 토지의 이용 상황을 실제 이용되고 있는 ‘자연림’으로 하여 개별공시지가를 산정한 다음 A감정평가법인에 검증을 의뢰하였는데, A감정평가법인이 그 토지의 이용 상황을 ‘공업용’으로 잘못 정정하여 검증지가를 산정하고, 시(市) 부동산평가위원회가 검증지가를 심의하면서 그 잘못을 발견하지 못하였다. 이에 따라 甲 소유 토지의 개별공시지가가 적정가격보다 훨씬 높은 가격으로 결정․공시되었다. B은행은 S시의 공시지가를 신뢰하고, 甲에게 70억원을 대출하였는데, 甲이 파산함에 따라 채권회수에 실패하였다. 다음 물음에 답하시오. (30점) \par
1) B은행은 S시를 대상으로 국가배상을 청구하였다. S시의 개별공시지가 결정행위가 국가배상법 제2조상의 위법행위에 해당하는가에 관하여 논하시오. (20점) \par
2) S시장은 개별공시지가제도의 입법목적을 이유로 S시 담당공무원들의 개별공시지가 산정에 관한 직무상 행위와 B은행의 손해사이에 상당인과관계가 없다고 항변한다. S시장의 항변의 타당성에 관하여 논하시오. (10점) \par
\end{longquestion}
